% Generated draft: requires mathematical and source review.
\documentclass[11pt]{article}
\usepackage[margin=25mm]{geometry}
\usepackage{fontspec}
\setmainfont{Latin Modern Roman}
\usepackage{amsmath,amssymb,mathtools,unicode-math}
\setmathfont{Latin Modern Math}
\usepackage{enumitem,needspace,xurl,hyperref}
\setlength{\parindent}{0pt}
\setlength{\parskip}{5pt}
\setlength{\emergencystretch}{3em}
\newcommand{\N}{\mathbb{N}}
\newcommand{\Z}{\mathbb{Z}}
\newcommand{\Q}{\mathbb{Q}}
\newcommand{\R}{\mathbb{R}}
\newcommand{\C}{\mathbb{C}}
\newcommand{\F}{\mathbb{F}}
\newcommand{\A}{\mathbb{A}}
\newcommand{\PP}{\mathbb{P}}
\newcommand{\E}{\mathbb{E}}
\newcommand{\eps}{\varepsilon}
\newcommand{\dd}{\,\mathrm{d}}
\begin{document}
\section*{UnsolvedMath ID: 11000207}
\textbf{Source problem code:} catalogue entry
\subsection{Definitions and mathematical statement}
Let \(G=\operatorname{Mod}(\Sigma)\) be the mapping class group in the source's surface convention. A discrete group is a-T-menable, equivalently has the Haagerup property, if it admits a metrically proper affine isometric action on a real Hilbert space. Concretely, seek a Hilbert space \(\mathcal H\), an orthogonal representation \(\pi:G\to O(\mathcal H)\) and a map \(b:G\to\mathcal H\) satisfying
\[b(gh)=b(g)+\pi(g)b(h),\]
so that \(g\cdot v=\pi(g)v+b(g)\), and
\[\forall R\geq0,\qquad\#\{g\in G:\lVert b(g)\rVert\leq R\}<\infty.\]
The question is whether such data exist for the specified mapping class group. The Hilbert space need not be finite-dimensional, and no cocompactness or free action is required. The condition is stronger than merely possessing an unbounded cocycle and is not a request for amenability.

The source excerpt omits the genus, puncture/boundary and orientation conventions; the research summary discusses \(g\geq2\), but a precise assertion must fix the surface, since low-complexity groups may have different behavior. The summary also calls absence of property (T) known without a supporting reference or exact surface range. That separate status assertion is not adopted as an established input, and even a verified failure of property (T) would not imply the existence of the required proper affine Hilbert action.
\subsection{Short English statement}
Does the intended mapping class group admit a proper affine isometric action on a Hilbert space? This is the Haagerup-property question, with surface conventions and the summary's separate property-(T) claim requiring verification.
\subsection{Sources}
\begin{itemize}
\item Local UnsolvedMath registry, permanent ID \(11000207\), exact source code \texttt{catalogue entry}.
\item Farb, \emph{Problems on Mapping Class Groups and Related Topics}, source heading ``Problem 8:'', PDF page 251; running chapter title \emph{Geometric properties of the mapping class group}: \url{https://www.math.uchicago.edu/~farb/papers/mcgbook.pdf}. Registry access 2026-07-29; pdf-font-structure extraction.
\item No precise primary reference supporting a current Haagerup or property-(T) status is supplied.
\end{itemize}
\end{document}
